
\section{模型评价}

本节对本文所建立的【N】类模型进行综合评价,分析各自的优点和局限性。

\subsection{【主方法1】模型}

\textbf{优点 1:}【说明】。

\textbf{优点 2:}【说明】。

\textbf{优点 3:}【说明】。

\textbf{优点 4:}【说明】。

\textbf{缺点 1:}【说明】。

\textbf{缺点 2:}【说明】。

\subsection{【主方法2】模型}

\textbf{优点 1:}【说明】。

\textbf{优点 2:}【说明】。

\textbf{优点 3:}【说明】。

\textbf{优点 4:}【说明】。

\textbf{缺点 1:}【说明】。

\textbf{缺点 2:}【说明】。

\subsection{【主方法3】模型}

\textbf{优点 1:}【说明】。

\textbf{优点 2:}【说明】。

\textbf{优点 3:}【说明】。

\textbf{优点 4:}【说明】。

\textbf{缺点 1:}【说明】。

\textbf{缺点 2:}【说明】。

\subsection{综合评价}

本文综合运用【方法1】、【方法2】、【方法3】和【方法4】四类方法,从不同角度对【研究问题】进行了系统建模。各类方法【优势互补/协同】,各问题的核心结论相互印证,形成了一个【研究链条】,为【应用场景】提供了有价值的参考。
